% !TeX root = ../navier-stokes-manuscript.tex
\subsection{Why finite-time escape opens the temporal Tower}\label{sub:BodyFiniteTimeResponseEscalation}

Suppose \(\norm{u(t)}_3\) escapes before \(T_*\). If \(u_t\) stayed uniformly bounded in \(L^3\), then
\[
  u(t)=u_0+\int_0^t u_t(s)\,ds
\]
would keep the velocity bounded on that finite interval. If \(u_{tt}\) stayed bounded, one integration would give
\[
  u_t(t)=u_t(0)+\int_0^t u_{tt}(s)\,ds
\]
and the first integral would again bound \(u\). Nothing changes at a higher fixed order: integrate down the derivative chain finitely many times and the velocity remains bounded.

This is why the singularity itself opens a temporal Tower. At fixed spatial coordinates, write
\begin{equation}\label{eq:TemporalTowerDefinitionBody}
  V_n:=\partial_t^nu,
  \qquad
  Q_n:=\partial_t^np,
  \qquad n\in\Nzero.
\end{equation}
The whole argument is the Banach-valued Taylor formula.

\begin{lemma}[Repeated finite-time integration]\label{lem:RepeatedFiniteTimeIntegration}
Let \(n\geq1\), and let \(f:[0,T_*)\to X\) have the derivatives required by the integral Taylor formula in a Banach space \(X\).
For \(0<t<T_*\),
\begin{equation}\label{eq:RepeatedFiniteTimeBound}
  \norm{f(t)}_X
  \leq
  \sum_{j=0}^{n-1}\frac{t^j}{j!}\norm{\partial_t^jf(0)}_X
  +\frac{t^n}{n!}\norm{\partial_t^nf}_{L^\infty(0,T_*;X)}.
\end{equation}
\end{lemma}

\begin{proof}
The integral Taylor formula gives
\[
  f(t)
  =\sum_{j=0}^{n-1}\frac{t^j}{j!}\partial_t^jf(0)
  +\frac1{(n-1)!}\int_0^t(t-s)^{n-1}\partial_t^nf(s)\,ds.
\]
Taking the \(X\)-norm and bounding the integral by its essential supremum gives \eqref{eq:RepeatedFiniteTimeBound}.
\end{proof}

Apply the lemma with \(f=u\) and \(X=L^3(\R^3)\). The smooth datum supplies every finite initial time derivative, and \(T_*<\infty\) keeps every polynomial factor finite. The escape in \eqref{eq:FiniteTimeLThreeEscape} can now happen only if, for each fixed \(n\),
\begin{equation}\label{eq:TemporalTowerMustEscape}
  \sup_{0<t<T_*}\norm{V_n(t)}_3=\infty.
\end{equation}
This conclusion is separate at every finite order. It gives no common blow-up point, rate, or sequence. It says exactly that no selected temporal response can stay bounded in the same critical space while the velocity escapes before a finite time. The Tower is
\[
  u,\quad u_t,\quad u_{tt},\quad u_{ttt},\quad\ldots.
\]
Acceleration and jerk help us feel the escalation, but a particle moving with the fluid does not experience \(u_t,u_{tt},\ldots\) alone. If \(X(a,t)\) is the trajectory beginning at the material label \(a\), then
\begin{equation}\label{eq:ParticleTrajectory}
  \frac d{dt}X(a,t)=u(X(a,t),t),
  \qquad
  X(a,0)=a,
\end{equation}
and
\begin{equation}\label{eq:MaterialAccelerationBody}
  \frac d{dt}u(X(a,t),t)
  =\partial_tu+(u\cdot\nabla)u
  =:D_tu.
\end{equation}
Along that trajectory, the equation gives
\begin{equation}\label{eq:MaterialAccelerationEquationBody}
  D_tu=-\nabla p+\nu\Delta u.
\end{equation}
One more material derivative differentiates pressure and curvature while also differentiating the moving frame through which the particle sees them. Time and space are already mixed in the physical response.

The fixed-coordinate rows let us keep that mixture in one frame. At order zero,
\begin{equation}\label{eq:BodyFirstTemporalVelocityRow}
  V_1+(V_0\cdot\nabla)V_0+\nabla Q_0
  =\nu\Delta V_0.
\end{equation}
Differentiating once gives
\begin{equation}\label{eq:BodySecondTemporalVelocityRow}
  V_2
  +(V_1\cdot\nabla)V_0
  +(V_0\cdot\nabla)V_1
  +\nabla Q_1
  =\nu\Delta V_1.
\end{equation}
The first nonlinear term changes the velocity doing the carrying. The second changes the response carried through the existing flow. On the right, viscosity acts two spatial derivatives above the current row. After \(n\) differentiations, the same division is
\begin{equation}\label{eq:BodyTemporalVelocityRecurrence}
  V_{n+1}
  +\sum_{a=0}^{n}\binom{n}{a}(V_a\cdot\nabla)V_{n-a}
  +\nabla Q_n
  =\nu\Delta V_n,
  \qquad
  \nabla\cdot V_n=0.
\end{equation}
Incompressibility fixes the pressure response from these same velocity rows:
\begin{equation}\label{eq:PressureNormalization}
  -\Delta p=\partial_i\partial_j(u_i u_j),
  \qquad
  p=R_iR_j(u_i u_j),
\end{equation}
with the decaying whole-space normalization.
Time differentiation gives
\begin{equation}\label{eq:BodyTemporalPressureRecurrence}
  Q_n
  =R_iR_j\sum_{a=0}^{n}\binom{n}{a}(V_a)_i(V_{n-a})_j.
\end{equation}
Every temporal step keeps all transport splits, reconstructs the pressure they create, and carries the viscous Laplacian with it. The singularity gave us the temporal direction. Spatial differentiation of these same rows gives the other direction.
